% --- Archon physics formalization source begin ---
% archon:physics
% archon:covers PhyXMiniProblems/problem_phyx_mini_0382.lean
% archon:source-report reports/phyx_mini/problem_phyx_mini_0382.source.json

\chapter{Physics problem phyx\_mini\_0382}
\label{ch:PhyXMiniProblems_problem_phyx_mini_0382}

\paragraph{Problem source.}
\#\# Physical scenario

Assume that the fluid is gasoline with atmospheric pressure 101 kPa on the top surface.

\#\# Question

What is the pressure at the bottom of the 7.5-m-tall storage tank of fluid at \$25\textasciicircum{}\{\textbackslash{}circ\}C\$ shown in figure?

\#\# Answer choices

- A: A: 146.2 kPa

- B: B: 156.2 kPa

- C: C: 256.2 kPa

- D: D: 166.2 kPa

\#\# Figure caption (auxiliary; use the image as primary evidence)

The image depicts a diagram of a spherical structure supported by several vertical columns. These columns are evenly spaced around the structure and are connected by diagonal braces for additional support. The entire assembly sits on a horizontal base. A vertical line with double-headed arrow indicates the height of the structure, labeled as \textbackslash{}( H \textbackslash{}). The sphere and supporting framework are illustrated in a simplified, schematic style.

\#\# Dataset metadata

Temperature and Heat Transfer; Physical Model Grounding Reasoning

\paragraph{Recorded answer/context.}
B: B: 156.2 kPa

\paragraph{Figure/image path.}
/root/proposal\_for\_physic/science-mango/phyx\_mini\_run/phyx\_data/test\_image/382.png

\paragraph{Formalization target.}
create a compiling Lean file with sorry bodies at `PhyXMiniProblems/problem\_phyx\_mini\_0382.lean`. The Lean declarations must preserve the physical quantities, dimensions or dimensional roles, figure labels, governing-law hypotheses, and final relation expressed by this problem.
Use Mathlib/Physlib names found through LeanExplore where available. If a physics API is missing, introduce faithful local abstractions rather than scalar placeholder aliases.

\begin{theorem}[Physics formalization target]
\label{thm:physics:phyx_mini_0382:target}
\lean{PhyXMiniProblems.ProblemPhyXMini0382.problem_phyx_mini_0382}
\uses{def:physics:phyx-mini-0382:phyxminiproblems-problemphyxmini0382-pressureinkilopascals, def:physics:phyx-mini-0382:phyxminiproblems-problemphyxmini0382-gasolinestoragetanksetup, def:physics:phyx-mini-0382:phyxminiproblems-problemphyxmini0382-matchesprimaryfigure, def:physics:phyx-mini-0382:phyxminiproblems-problemphyxmini0382-matchesproblemstatement, def:physics:phyx-mini-0382:phyxminiproblems-problemphyxmini0382-usesreferencegasolineandearthgravity, def:physics:phyx-mini-0382:phyxminiproblems-problemphyxmini0382-hasphysicaltankparameters, def:physics:phyx-mini-0382:phyxminiproblems-problemphyxmini0382-satisfieshydrostaticpressurelaw, def:physics:phyx-mini-0382:phyxminiproblems-problemphyxmini0382-matchesdisplayedbottompressure, def:physics:phyx-mini-0382:phyxminiproblems-problemphyxmini0382-isuniquematchingdisplayedpressure}
The assigned autoformalize agent should translate this physics problem into Lean declarations in the covered file, with theorem and lemma proof bodies written as `by sorry`.
\end{theorem}
\begin{proof}
This is an autoformalization task, not a proof task. Produce statements that can later be proved by the physics prover without weakening the physical model.
\end{proof}

% --- Archon declaration topology begin ---
\section{Declaration topology}
\begin{definition}[Length Quantity]
  \label{def:physics:phyx-mini-0382:phyxminiproblems-problemphyxmini0382-lengthquantity}
  \lean{PhyXMiniProblems.ProblemPhyXMini0382.LengthQuantity}
  A nonnegative physical length, independent of the chosen readout unit
\end{definition}
\begin{proof}
  This is the declared model definition of Length Quantity.
\end{proof}
\begin{definition}[Mass Density Quantity]
  \label{def:physics:phyx-mini-0382:phyxminiproblems-problemphyxmini0382-massdensityquantity}
  \lean{PhyXMiniProblems.ProblemPhyXMini0382.MassDensityQuantity}
  A nonnegative physical mass density, with dimension M L⁻³
\end{definition}
\begin{proof}
  This is the declared model definition of Mass Density Quantity.
\end{proof}
\begin{definition}[Acceleration Quantity]
  \label{def:physics:phyx-mini-0382:phyxminiproblems-problemphyxmini0382-accelerationquantity}
  \lean{PhyXMiniProblems.ProblemPhyXMini0382.AccelerationQuantity}
  A nonnegative acceleration magnitude, with dimension L T⁻²
\end{definition}
\begin{proof}
  This is the declared model definition of Acceleration Quantity.
\end{proof}
\begin{definition}[Pressure Quantity]
  \label{def:physics:phyx-mini-0382:phyxminiproblems-problemphyxmini0382-pressurequantity}
  \lean{PhyXMiniProblems.ProblemPhyXMini0382.PressureQuantity}
  Absolute pressure, represented by Physlib's dimensional pressure type
\end{definition}
\begin{proof}
  This is the declared model definition of Pressure Quantity.
\end{proof}
\begin{definition}[Temperature Quantity]
  \label{def:physics:phyx-mini-0382:phyxminiproblems-problemphyxmini0382-temperaturequantity}
  \lean{PhyXMiniProblems.ProblemPhyXMini0382.TemperatureQuantity}
  Absolute thermodynamic temperature
\end{definition}
\begin{proof}
  This is the declared model definition of Temperature Quantity.
\end{proof}
\begin{definition}[length Readout]
  \label{def:physics:phyx-mini-0382:phyxminiproblems-problemphyxmini0382-lengthreadout}
  \lean{PhyXMiniProblems.ProblemPhyXMini0382.lengthReadout}
  \uses{def:physics:phyx-mini-0382:phyxminiproblems-problemphyxmini0382-lengthquantity}
  Read a physical length in a selected length unit
\end{definition}
\begin{proof}
  This is the declared model definition of length Readout.
\end{proof}
\begin{definition}[pressure Readout]
  \label{def:physics:phyx-mini-0382:phyxminiproblems-problemphyxmini0382-pressurereadout}
  \lean{PhyXMiniProblems.ProblemPhyXMini0382.pressureReadout}
  \uses{def:physics:phyx-mini-0382:phyxminiproblems-problemphyxmini0382-pressurequantity}
  Read pressure in the coherent unit induced by selected base units
\end{definition}
\begin{proof}
  This is the declared model definition of pressure Readout.
\end{proof}
\begin{definition}[density Readout]
  \label{def:physics:phyx-mini-0382:phyxminiproblems-problemphyxmini0382-densityreadout}
  \lean{PhyXMiniProblems.ProblemPhyXMini0382.densityReadout}
  \uses{def:physics:phyx-mini-0382:phyxminiproblems-problemphyxmini0382-massdensityquantity}
  Read mass density in a selected mass unit per selected length unit cubed
\end{definition}
\begin{proof}
  This is the declared model definition of density Readout.
\end{proof}
\begin{definition}[acceleration Readout]
  \label{def:physics:phyx-mini-0382:phyxminiproblems-problemphyxmini0382-accelerationreadout}
  \lean{PhyXMiniProblems.ProblemPhyXMini0382.accelerationReadout}
  \uses{def:physics:phyx-mini-0382:phyxminiproblems-problemphyxmini0382-accelerationquantity}
  Read acceleration in selected length and time units
\end{definition}
\begin{proof}
  This is the declared model definition of acceleration Readout.
\end{proof}
\begin{definition}[length In Meters]
  \label{def:physics:phyx-mini-0382:phyxminiproblems-problemphyxmini0382-lengthinmeters}
  \lean{PhyXMiniProblems.ProblemPhyXMini0382.lengthInMeters}
  \uses{def:physics:phyx-mini-0382:phyxminiproblems-problemphyxmini0382-lengthquantity, def:physics:phyx-mini-0382:phyxminiproblems-problemphyxmini0382-lengthreadout}
  Meter readout of a physical length
\end{definition}
\begin{proof}
  This is the declared model definition of length In Meters.
\end{proof}
\begin{definition}[pressure In Pascals]
  \label{def:physics:phyx-mini-0382:phyxminiproblems-problemphyxmini0382-pressureinpascals}
  \lean{PhyXMiniProblems.ProblemPhyXMini0382.pressureInPascals}
  \uses{def:physics:phyx-mini-0382:phyxminiproblems-problemphyxmini0382-pressurequantity, def:physics:phyx-mini-0382:phyxminiproblems-problemphyxmini0382-pressurereadout}
  Pascal readout of an absolute pressure
\end{definition}
\begin{proof}
  This is the declared model definition of pressure In Pascals.
\end{proof}
\begin{definition}[pressure In Kilopascals]
  \label{def:physics:phyx-mini-0382:phyxminiproblems-problemphyxmini0382-pressureinkilopascals}
  \lean{PhyXMiniProblems.ProblemPhyXMini0382.pressureInKilopascals}
  \uses{def:physics:phyx-mini-0382:phyxminiproblems-problemphyxmini0382-pressurequantity, def:physics:phyx-mini-0382:phyxminiproblems-problemphyxmini0382-pressureinpascals}
  Kilopascal readout of an absolute pressure
\end{definition}
\begin{proof}
  This is the declared model definition of pressure In Kilopascals.
\end{proof}
\begin{definition}[density In Kilograms Per Cubic Meter]
  \label{def:physics:phyx-mini-0382:phyxminiproblems-problemphyxmini0382-densityinkilogramspercubicmeter}
  \lean{PhyXMiniProblems.ProblemPhyXMini0382.densityInKilogramsPerCubicMeter}
  \uses{def:physics:phyx-mini-0382:phyxminiproblems-problemphyxmini0382-massdensityquantity, def:physics:phyx-mini-0382:phyxminiproblems-problemphyxmini0382-densityreadout}
  Kilogram-per-cubic-meter readout of a physical mass density
\end{definition}
\begin{proof}
  This is the declared model definition of density In Kilograms Per Cubic Meter.
\end{proof}
\begin{definition}[acceleration In Meters Per Second Squared]
  \label{def:physics:phyx-mini-0382:phyxminiproblems-problemphyxmini0382-accelerationinmeterspersecondsquared}
  \lean{PhyXMiniProblems.ProblemPhyXMini0382.accelerationInMetersPerSecondSquared}
  \uses{def:physics:phyx-mini-0382:phyxminiproblems-problemphyxmini0382-accelerationquantity, def:physics:phyx-mini-0382:phyxminiproblems-problemphyxmini0382-accelerationreadout}
  Meter-per-second-squared readout of an acceleration magnitude
\end{definition}
\begin{proof}
  This is the declared model definition of acceleration In Meters Per Second Squared.
\end{proof}
\begin{definition}[temperature In Kelvins]
  \label{def:physics:phyx-mini-0382:phyxminiproblems-problemphyxmini0382-temperatureinkelvins}
  \lean{PhyXMiniProblems.ProblemPhyXMini0382.temperatureInKelvins}
  \uses{def:physics:phyx-mini-0382:phyxminiproblems-problemphyxmini0382-temperaturequantity}
  Read a stored absolute temperature in kelvins
\end{definition}
\begin{proof}
  This is the declared model definition of temperature In Kelvins.
\end{proof}
\begin{definition}[temperature In Degrees Celsius]
  \label{def:physics:phyx-mini-0382:phyxminiproblems-problemphyxmini0382-temperatureindegreescelsius}
  \lean{PhyXMiniProblems.ProblemPhyXMini0382.temperatureInDegreesCelsius}
  \uses{def:physics:phyx-mini-0382:phyxminiproblems-problemphyxmini0382-temperaturequantity, def:physics:phyx-mini-0382:phyxminiproblems-problemphyxmini0382-temperatureinkelvins}
  Celsius readout, using the exact offset 0 °C = 273.15 K
\end{definition}
\begin{proof}
  This is the declared model definition of temperature In Degrees Celsius.
\end{proof}
\begin{definition}[Tank Shape]
  \label{def:physics:phyx-mini-0382:phyxminiproblems-problemphyxmini0382-tankshape}
  \lean{PhyXMiniProblems.ProblemPhyXMini0382.TankShape}
  Shape of the storage vessel visible in the primary bitmap
\end{definition}
\begin{proof}
  This is the declared model definition of Tank Shape.
\end{proof}
\begin{definition}[Tank Liquid]
  \label{def:physics:phyx-mini-0382:phyxminiproblems-problemphyxmini0382-tankliquid}
  \lean{PhyXMiniProblems.ProblemPhyXMini0382.TankLiquid}
  Liquid named in the problem statement
\end{definition}
\begin{proof}
  This is the declared model definition of Tank Liquid.
\end{proof}
\begin{definition}[Top Surface Boundary]
  \label{def:physics:phyx-mini-0382:phyxminiproblems-problemphyxmini0382-topsurfaceboundary}
  \lean{PhyXMiniProblems.ProblemPhyXMini0382.TopSurfaceBoundary}
  Boundary condition at the upper gasoline surface
\end{definition}
\begin{proof}
  This is the declared model definition of Top Surface Boundary.
\end{proof}
\begin{definition}[Tank Thermal Condition]
  \label{def:physics:phyx-mini-0382:phyxminiproblems-problemphyxmini0382-tankthermalcondition}
  \lean{PhyXMiniProblems.ProblemPhyXMini0382.TankThermalCondition}
  Thermal condition used for the tabulated gasoline density
\end{definition}
\begin{proof}
  This is the declared model definition of Tank Thermal Condition.
\end{proof}
\begin{definition}[Figure Quantity Label]
  \label{def:physics:phyx-mini-0382:phyxminiproblems-problemphyxmini0382-figurequantitylabel}
  \lean{PhyXMiniProblems.ProblemPhyXMini0382.FigureQuantityLabel}
  The literal quantity label visible beside the height arrow
\end{definition}
\begin{proof}
  This is the declared model definition of Figure Quantity Label.
\end{proof}
\begin{definition}[Figure Height Role]
  \label{def:physics:phyx-mini-0382:phyxminiproblems-problemphyxmini0382-figureheightrole}
  \lean{PhyXMiniProblems.ProblemPhyXMini0382.FigureHeightRole}
  Physical role assigned to the figure's H marker
\end{definition}
\begin{proof}
  This is the declared model definition of Figure Height Role.
\end{proof}
\begin{definition}[Figure Height Endpoint]
  \label{def:physics:phyx-mini-0382:phyxminiproblems-problemphyxmini0382-figureheightendpoint}
  \lean{PhyXMiniProblems.ProblemPhyXMini0382.FigureHeightEndpoint}
  Endpoints of the vertical double-headed arrow in the supplied figure
\end{definition}
\begin{proof}
  This is the declared model definition of Figure Height Endpoint.
\end{proof}
\begin{definition}[Spherical Storage Tank Figure]
  \label{def:physics:phyx-mini-0382:phyxminiproblems-problemphyxmini0382-sphericalstoragetankfigure}
  \lean{PhyXMiniProblems.ProblemPhyXMini0382.SphericalStorageTankFigure}
  \uses{def:physics:phyx-mini-0382:phyxminiproblems-problemphyxmini0382-lengthquantity, def:physics:phyx-mini-0382:phyxminiproblems-problemphyxmini0382-tankshape, def:physics:phyx-mini-0382:phyxminiproblems-problemphyxmini0382-figurequantitylabel, def:physics:phyx-mini-0382:phyxminiproblems-problemphyxmini0382-figureheightrole, def:physics:phyx-mini-0382:phyxminiproblems-problemphyxmini0382-figureheightendpoint}
  Structured transcription of image 382.png
\end{definition}
\begin{proof}
  This is the declared model definition of Spherical Storage Tank Figure.
\end{proof}
\begin{definition}[Gasoline Storage Tank Setup]
  \label{def:physics:phyx-mini-0382:phyxminiproblems-problemphyxmini0382-gasolinestoragetanksetup}
  \lean{PhyXMiniProblems.ProblemPhyXMini0382.GasolineStorageTankSetup}
  \uses{def:physics:phyx-mini-0382:phyxminiproblems-problemphyxmini0382-lengthquantity, def:physics:phyx-mini-0382:phyxminiproblems-problemphyxmini0382-massdensityquantity, def:physics:phyx-mini-0382:phyxminiproblems-problemphyxmini0382-accelerationquantity, def:physics:phyx-mini-0382:phyxminiproblems-problemphyxmini0382-pressurequantity, def:physics:phyx-mini-0382:phyxminiproblems-problemphyxmini0382-temperaturequantity, def:physics:phyx-mini-0382:phyxminiproblems-problemphyxmini0382-tankliquid, def:physics:phyx-mini-0382:phyxminiproblems-problemphyxmini0382-topsurfaceboundary, def:physics:phyx-mini-0382:phyxminiproblems-problemphyxmini0382-tankthermalcondition, def:physics:phyx-mini-0382:phyxminiproblems-problemphyxmini0382-sphericalstoragetankfigure}
  Independent physical quantities in the tank problem. In particular, bottomAbsolutePressure is an observable field rather than a definition made from the recorded answer or the hydrostatic formula
\end{definition}
\begin{proof}
  This is the declared model definition of Gasoline Storage Tank Setup.
\end{proof}
\begin{definition}[Matches Primary Figure]
  \label{def:physics:phyx-mini-0382:phyxminiproblems-problemphyxmini0382-matchesprimaryfigure}
  \lean{PhyXMiniProblems.ProblemPhyXMini0382.MatchesPrimaryFigure}
  \uses{def:physics:phyx-mini-0382:phyxminiproblems-problemphyxmini0382-gasolinestoragetanksetup}
  Qualitative and labelled geometry read directly from the primary bitmap. The figure itself supplies no numerical value for H
\end{definition}
\begin{proof}
  This is the declared model definition of Matches Primary Figure.
\end{proof}
\begin{definition}[Matches Problem Statement]
  \label{def:physics:phyx-mini-0382:phyxminiproblems-problemphyxmini0382-matchesproblemstatement}
  \lean{PhyXMiniProblems.ProblemPhyXMini0382.MatchesProblemStatement}
  \uses{def:physics:phyx-mini-0382:phyxminiproblems-problemphyxmini0382-lengthinmeters, def:physics:phyx-mini-0382:phyxminiproblems-problemphyxmini0382-pressureinkilopascals, def:physics:phyx-mini-0382:phyxminiproblems-problemphyxmini0382-temperatureindegreescelsius, def:physics:phyx-mini-0382:phyxminiproblems-problemphyxmini0382-gasolinestoragetanksetup}
  Quantitative and qualitative information stated in the problem. The fluid is filled to its atmospheric top surface, so its hydrostatic depth is the stated 7.5 m tank height. No bottom-pressure value occurs here
\end{definition}
\begin{proof}
  This is the declared model definition of Matches Problem Statement.
\end{proof}
\begin{definition}[Uses Reference Gasoline And Earth Gravity]
  \label{def:physics:phyx-mini-0382:phyxminiproblems-problemphyxmini0382-usesreferencegasolineandearthgravity}
  \lean{PhyXMiniProblems.ProblemPhyXMini0382.UsesReferenceGasolineAndEarthGravity}
  \uses{def:physics:phyx-mini-0382:phyxminiproblems-problemphyxmini0382-densityinkilogramspercubicmeter, def:physics:phyx-mini-0382:phyxminiproblems-problemphyxmini0382-accelerationinmeterspersecondsquared, def:physics:phyx-mini-0382:phyxminiproblems-problemphyxmini0382-gasolinestoragetanksetup}
  Standard reference-property values implicit in the recorded multiple-choice answer: gasoline at 25 °C has density 750 kg/m³, and near-surface Earth gravity is 9.81 m/s². These calibrations do not mention bottom pressure
\end{definition}
\begin{proof}
  This is the declared model definition of Uses Reference Gasoline And Earth Gravity.
\end{proof}
\begin{definition}[Has Physical Tank Parameters]
  \label{def:physics:phyx-mini-0382:phyxminiproblems-problemphyxmini0382-hasphysicaltankparameters}
  \lean{PhyXMiniProblems.ProblemPhyXMini0382.HasPhysicalTankParameters}
  \uses{def:physics:phyx-mini-0382:phyxminiproblems-problemphyxmini0382-lengthinmeters, def:physics:phyx-mini-0382:phyxminiproblems-problemphyxmini0382-pressureinpascals, def:physics:phyx-mini-0382:phyxminiproblems-problemphyxmini0382-densityinkilogramspercubicmeter, def:physics:phyx-mini-0382:phyxminiproblems-problemphyxmini0382-accelerationinmeterspersecondsquared, def:physics:phyx-mini-0382:phyxminiproblems-problemphyxmini0382-gasolinestoragetanksetup}
  Positivity of the physical parameters and absolute pressures
\end{definition}
\begin{proof}
  This is the declared model definition of Has Physical Tank Parameters.
\end{proof}
\begin{definition}[Satisfies Hydrostatic Pressure Law]
  \label{def:physics:phyx-mini-0382:phyxminiproblems-problemphyxmini0382-satisfieshydrostaticpressurelaw}
  \lean{PhyXMiniProblems.ProblemPhyXMini0382.SatisfiesHydrostaticPressureLaw}
  \uses{def:physics:phyx-mini-0382:phyxminiproblems-problemphyxmini0382-lengthreadout, def:physics:phyx-mini-0382:phyxminiproblems-problemphyxmini0382-pressurereadout, def:physics:phyx-mini-0382:phyxminiproblems-problemphyxmini0382-densityreadout, def:physics:phyx-mini-0382:phyxminiproblems-problemphyxmini0382-accelerationreadout, def:physics:phyx-mini-0382:phyxminiproblems-problemphyxmini0382-gasolinestoragetanksetup}
  For a stationary uniform-density liquid column, the absolute pressure gain from the free surface to the bottom is ρ g h. The equation is stated for every coherent choice of mass, length, and time units. It contains no numerical bottom-pressure answer
\end{definition}
\begin{proof}
  This is the declared model definition of Satisfies Hydrostatic Pressure Law.
\end{proof}
\begin{definition}[Answer Choice]
  \label{def:physics:phyx-mini-0382:phyxminiproblems-problemphyxmini0382-answerchoice}
  \lean{PhyXMiniProblems.ProblemPhyXMini0382.AnswerChoice}
  Labels of the four displayed answer choices
\end{definition}
\begin{proof}
  This is the declared model definition of Answer Choice.
\end{proof}
\begin{definition}[displayed Pressure Kilopascals]
  \label{def:physics:phyx-mini-0382:phyxminiproblems-problemphyxmini0382-displayedpressurekilopascals}
  \lean{PhyXMiniProblems.ProblemPhyXMini0382.displayedPressureKilopascals}
  \uses{def:physics:phyx-mini-0382:phyxminiproblems-problemphyxmini0382-answerchoice}
  Kilopascal value printed beside each answer label
\end{definition}
\begin{proof}
  This is the declared model definition of displayed Pressure Kilopascals.
\end{proof}
\begin{definition}[recorded Dataset Answer]
  \label{def:physics:phyx-mini-0382:phyxminiproblems-problemphyxmini0382-recordeddatasetanswer}
  \lean{PhyXMiniProblems.ProblemPhyXMini0382.recordedDatasetAnswer}
  \uses{def:physics:phyx-mini-0382:phyxminiproblems-problemphyxmini0382-answerchoice}
  Dataset answer metadata, deliberately not used as a premise
\end{definition}
\begin{proof}
  This is the declared model definition of recorded Dataset Answer.
\end{proof}
\begin{definition}[Rounds To Nearest Tenth Kilopascal]
  \label{def:physics:phyx-mini-0382:phyxminiproblems-problemphyxmini0382-roundstonearesttenthkilopascal}
  \lean{PhyXMiniProblems.ProblemPhyXMini0382.RoundsToNearestTenthKilopascal}
  \uses{def:physics:phyx-mini-0382:phyxminiproblems-problemphyxmini0382-pressurequantity, def:physics:phyx-mini-0382:phyxminiproblems-problemphyxmini0382-pressureinkilopascals}
  A physical pressure rounds to a displayed tenth of a kilopascal
\end{definition}
\begin{proof}
  This is the declared model definition of Rounds To Nearest Tenth Kilopascal.
\end{proof}
\begin{definition}[Matches Displayed Bottom Pressure]
  \label{def:physics:phyx-mini-0382:phyxminiproblems-problemphyxmini0382-matchesdisplayedbottompressure}
  \lean{PhyXMiniProblems.ProblemPhyXMini0382.MatchesDisplayedBottomPressure}
  \uses{def:physics:phyx-mini-0382:phyxminiproblems-problemphyxmini0382-gasolinestoragetanksetup, def:physics:phyx-mini-0382:phyxminiproblems-problemphyxmini0382-answerchoice, def:physics:phyx-mini-0382:phyxminiproblems-problemphyxmini0382-displayedpressurekilopascals, def:physics:phyx-mini-0382:phyxminiproblems-problemphyxmini0382-roundstonearesttenthkilopascal}
  A displayed choice agrees with the modeled bottom pressure
\end{definition}
\begin{proof}
  This is the declared model definition of Matches Displayed Bottom Pressure.
\end{proof}
\begin{definition}[Is Unique Matching Displayed Pressure]
  \label{def:physics:phyx-mini-0382:phyxminiproblems-problemphyxmini0382-isuniquematchingdisplayedpressure}
  \lean{PhyXMiniProblems.ProblemPhyXMini0382.IsUniqueMatchingDisplayedPressure}
  \uses{def:physics:phyx-mini-0382:phyxminiproblems-problemphyxmini0382-gasolinestoragetanksetup, def:physics:phyx-mini-0382:phyxminiproblems-problemphyxmini0382-answerchoice, def:physics:phyx-mini-0382:phyxminiproblems-problemphyxmini0382-matchesdisplayedbottompressure}
  Exactly one displayed answer agrees with the modeled bottom pressure
\end{definition}
\begin{proof}
  This is the declared model definition of Is Unique Matching Displayed Pressure.
\end{proof}
\begin{lemma}[bottom pressure in kilopascals]
  \label{lem:physics:phyx-mini-0382:phyxminiproblems-problemphyxmini0382-bottom-pressure-in-kilopascals}
  \lean{PhyXMiniProblems.ProblemPhyXMini0382.bottom_pressure_in_kilopascals}
  \uses{def:physics:phyx-mini-0382:phyxminiproblems-problemphyxmini0382-pressureinkilopascals, def:physics:phyx-mini-0382:phyxminiproblems-problemphyxmini0382-gasolinestoragetanksetup, def:physics:phyx-mini-0382:phyxminiproblems-problemphyxmini0382-matchesproblemstatement, def:physics:phyx-mini-0382:phyxminiproblems-problemphyxmini0382-usesreferencegasolineandearthgravity, def:physics:phyx-mini-0382:phyxminiproblems-problemphyxmini0382-satisfieshydrostaticpressurelaw}
  The hydrostatic law and the calibrated data give the unrounded bottom pressure 156.18125 kPa = 124945 / 800 kPa. This derived value is not a premise
\end{lemma}
\begin{proof}
  The conclusion follows from the governing relations and figure readouts named in the statement of bottom pressure in kilopascals.
\end{proof}
% --- Archon declaration topology end ---
% --- Archon physics formalization source end ---
